\begin{abstract}
Micro, Small, and Medium Enterprises (MSMEs) make up the largest segment of businesses worldwide but experience high failure rates due to poor access to finance and lack of reliable viability assessment techniques. Previous machine learning models for MSME and small business evaluation rely solely on binary categorization (default vs.\ non-default), provide no risk stratification, and deliver no advice for enhancing the success of enterprises. We introduce an AI-based solution that overcomes these limitations in three key ways: (1)~a unique five-class viability taxonomy (Critical, At-Risk, Stable, Growing, Thriving) that is generated using a health score, which converts loan outcomes from binary classifications into risk gradations; (2)~per-prediction interpretability using TreeSHAP; and (3)~a prescriptive counterfactual model that recommends minimal changes that will improve the classification of viability. By training XGBoost and LightGBM on 899,164 U.S.\ SBA loan transactions, each with 11 engineered variables, we attain an average accuracy rate of 92.78\% and 92.98\%, respectively, in five-class predictions (macro AUC $>$ 0.986). In contrast, logistic regression attains an average accuracy of only 77.67\%, with naive baselines scoring 53.31\%. Moreover, the counterfactual engine successfully generates viable solutions for 200 examples at sub-3-millisecond latency rates, resolving 82\% of transitions by altering a single feature value. SHAP and counterfactual global analysis are highly internally consistent, with loan tenure and gross approved amount serving as the key viability indicators in both models. This system has been implemented as a FastAPI--Streamlit decoupled architecture with auditing capabilities, providing proof that viability prediction, explainability, and prescriptive action can coexist in one system.
\end{abstract}

\keywords{MSME viability assessment \and XGBoost \and LightGBM \and SHAP explainability \and counterfactual recommendations \and multi-class classification \and credit risk \and prescriptive analytics}
